\section{The Value of a Distinction}
\label{sec:value}

Minimal refinement says how to add a distinction; it does not license every possible distinction. The value question is task-dependent: when does the new distinction actually change the best achievable prediction? We measure this by the decrease in optimal risk after the information structure is refined, in the same decision-theoretic sense used in classical comparisons of statistical experiments and value of information \citep{blackwell1953equivalent,berger1985statistical}.

\subsection{Risk monotonicity}

The first fact is purely order-theoretic. If $\calG$ contains all events in $\calF$, then every $\calF$-measurable predictor is also $\calG$-measurable. The decision class has only grown; this is the Bayes-risk analogue of Blackwell's monotonicity of decision value under more informative experiments \citep{blackwell1953equivalent}.

\begin{theorem}[Bayes-risk monotonicity]
\label{thm:risk-monotonicity:informal}
If $\calF\subseteq\calG$, then for every loss $L$ for which both risks are well-defined,
\begin{align*}
R_L(\calG)\le R_L(\calF).
\end{align*}
\end{theorem}

This is a population statement. It says that information cannot hurt the best possible decision rule. It does not say that a finite implementation should refine without restraint. Statistical estimation, finite samples, and model selection can overfit; those are algorithmic questions, not failures of the information order. This is why the DGM algorithm later uses penalties, thresholds, and sample-split scoring before committing a new concept.

\subsection{Log loss}

For a finite target variable $Y$, log loss asks the learner to output a conditional distribution. The optimal prediction under $\calF$ is the conditional law of $Y$ given $\calF$, and the optimal risk is conditional entropy, a standard identity for logarithmic scoring rules and information theory \citep{cover2006elements,gneiting2007strictly}.

\begin{theorem}[Log-loss value]
\label{thm:log-value:informal}
Let $Y$ be finite and let the loss be $-\log q(Y)$. If $\calF\subseteq\calG$, then
\begin{align*}
R_{\log}(\calF)-R_{\log}(\calG)=I(Y;\calG\mid\calF).
\end{align*}
\end{theorem}

Here $I(Y;\calG\mid\calF)$ denotes $H(Y\mid\calF)-H(Y\mid\calG)$. Thus, for classification and language modeling, a refinement is useful exactly to the extent that it carries conditional mutual information about the consequence beyond what the learner already knows. If the new bit does not tell us anything about the next consequence beyond the old information, it is a gratuitous distinction, even if it separates inputs.

\subsection{Squared loss}

For $Y\in L^2(\Omega;\R^d)$ and squared loss, the optimal predictor under $\calF$ is $\E[Y\mid\calF]$, the $L^2$ projection onto the closed subspace of $\calF$-measurable random variables \citep{billingsley1995probability,kallenberg2002foundations}. Refinement helps precisely when it changes the conditional mean.

\begin{theorem}[Squared-loss value]
\label{thm:squared-value:informal}
Let $Y\in L^2(\Omega;\R^d)$ and let $L(Y,a)=\norm{Y-a}_2^2$. If $\calF\subseteq\calG$, then
\begin{align*}
R_2(\calF)-R_2(\calG)=\E\left[\norm{\E[Y\mid\calG]-\E[Y\mid\calF]}_2^2\right].
\end{align*}
\end{theorem}

The identity is the Hilbert-space analogue of the log-loss result. A distinction has squared-loss value when the cells it creates have different conditional means. In regression language, refinement is worth paying for when it separates consequences, not merely when it separates inputs.

\subsection{Finite empirical scores}

The population identities lead directly to standard empirical split scores: information gain for classification and variance reduction for regression trees \citep{breiman1984classification,domingos2000mining}. Suppose an atom $B$ contains samples $(x_j,y_j)$ and a candidate distinction $g:B\to\{0,1\}$ splits it into $B_0$ and $B_1$.

For classification, let $\hat p_B$, $\hat p_{B_0}$, and $\hat p_{B_1}$ be empirical label distributions. The empirical log-loss improvement is
\begin{align*}
\Delta_{\log}(B,g)=\abs{B}H(\hat p_B)-\abs{B_0}H(\hat p_{B_0})-\abs{B_1}H(\hat p_{B_1}).
\end{align*}
Equivalently,
\begin{align*}
\Delta_{\log}(B,g)=\abs{B}\widehat I(Y;g\mid B).
\end{align*}

For regression, define
\begin{align*}
\hat\mu_B=\frac{1}{\abs{B}}\sum_{j\in B}y_j,\qquad \operatorname{SSE}(B)=\sum_{j\in B}\norm{y_j-\hat\mu_B}_2^2.
\end{align*}
The empirical squared-loss improvement is
\begin{align*}
\Delta_2(B,g)=\operatorname{SSE}(B)-\operatorname{SSE}(B_0)-\operatorname{SSE}(B_1).
\end{align*}
It also has the between-cell form
\begin{align*}
\Delta_2(B,g)=\abs{B_0}\norm{\hat\mu_{B_0}-\hat\mu_B}_2^2+\abs{B_1}\norm{\hat\mu_{B_1}-\hat\mu_B}_2^2.
\end{align*}

These scores implement the same principle as the population theory: split only when the proposed distinction has measurable consequence value. The DGM algorithm below turns this population rule into a guarded empirical decision.
